\section{Evaluate Boolean Expression to True}
\label{sec:dp-evaluate-boolean-expression}

\problemstatement{Given a boolean expression string containing
operands \texttt{T} and \texttt{F} separated by operators
\texttt{\&} (and), \texttt{|} (or), and \texttt{\^{}} (xor), count the
number of ways to fully parenthesize the expression so that it
evaluates to \texttt{True} (answer modulo $10^9+7$).}

\begin{patternbox}
\emph{``Count the parenthesizations achieving a target outcome''} is
this chapter's interval template with two simultaneous twists: the
split points $k$ are restricted to \textbf{operator positions} only
(operands themselves cannot be split), and each range tracks a
\textbf{pair} of counts (ways \texttt{True}, ways \texttt{False})
rather than one number -- because combining two ranges' outcomes through
an operator needs to know \emph{both} sub-outcomes' counts to compute
the combined result correctly.
\end{patternbox}

\subsection*{Deriving the combination rules}

Let $f(i,j) = (\textit{trueWays}, \textit{falseWays})$ for the
sub-expression spanning operand positions $i$ to $j$. Base case:
$i=j$ (a single operand): $f(i,i) = (1,0)$ if $\textit{expr}[i]=
\texttt{T}$, else $(0,1)$. For $i<j$, try every operator position $k$
strictly between $i$ and $j$, splitting into left $(i,k{-}1)$ and right
$(k{+}1,j)$, and combine according to which operator sits at $k$:

\begin{itemize}
  \item \texttt{\&} (and): result is \texttt{True} only if
        \emph{both} sides are \texttt{True}.
        $\textit{trueWays} \mathrel{+}= L_T \cdot R_T$.
        $\textit{falseWays} \mathrel{+}= L_T R_F + L_F R_T + L_F R_F$.
  \item \texttt{|} (or): result is \texttt{False} only if \emph{both}
        sides are \texttt{False}.
        $\textit{trueWays} \mathrel{+}= L_T R_T + L_T R_F + L_F R_T$.
        $\textit{falseWays} \mathrel{+}= L_F R_F$.
  \item \texttt{\^{}} (xor): result is \texttt{True} only if the sides
        \emph{differ}.
        $\textit{trueWays} \mathrel{+}= L_T R_F + L_F R_T$.
        $\textit{falseWays} \mathrel{+}= L_T R_T + L_F R_F$.
\end{itemize}
Sum these contributions over every valid operator position $k$.

\subsection*{C++ implementation (memoization)}

\begin{lstlisting}
#include <bits/stdc++.h>
using namespace std;
const int MOD = 1e9 + 7;

pair<long long, long long> evalMemo(int i, int j, string& expr,
                                     vector<vector<pair<long long,long long>>>& memo) {
    if (i == j) {
        return expr[i] == 'T' ? make_pair(1LL, 0LL) : make_pair(0LL, 1LL);
    }
    if (memo[i][j].first != -1) return memo[i][j];

    long long trueWays = 0, falseWays = 0;
    for (int k = i + 1; k < j; k += 2) { // operators sit at odd offsets
        auto [lt, lf] = evalMemo(i, k - 1, expr, memo);
        auto [rt, rf] = evalMemo(k + 1, j, expr, memo);

        if (expr[k] == '&') {
            trueWays = (trueWays + lt * rt) % MOD;
            falseWays = (falseWays + lt * rf + lf * rt + lf * rf) % MOD;
        } else if (expr[k] == '|') {
            trueWays = (trueWays + lt * rt + lt * rf + lf * rt) % MOD;
            falseWays = (falseWays + lf * rf) % MOD;
        } else { // '^'
            trueWays = (trueWays + lt * rf + lf * rt) % MOD;
            falseWays = (falseWays + lt * rt + lf * rf) % MOD;
        }
    }
    return memo[i][j] = {trueWays, falseWays};
}

int evaluateBooleanExpr(string expr) {
    int n = expr.size();
    vector<vector<pair<long long,long long>>> memo(n, vector<pair<long long,long long>>(n, {-1, -1}));
    return evalMemo(0, n - 1, expr, memo).first;
}

int main() {
    cout << evaluateBooleanExpr("T|F&T^F") << endl; // example expression
    return 0;
}
\end{lstlisting}

\subsection*{Dry run}

\dryrun{Take $\textit{expr} = \texttt{"T\^{}F|F"}$ (positions:
$0{=}\texttt{T}$, $1{=}\texttt{\^{}}$, $2{=}\texttt{F}$,
$3{=}\texttt{|}$, $4{=}\texttt{F}$).}

Base cases: $f(0,0)=(1,0)$, $f(2,2)=(0,1)$, $f(4,4)=(0,1)$.
$f(0,2)$ (sub-expression \texttt{"T\^{}F"}, operator \texttt{\^{}} at
position $1$): $L=(1,0)$, $R=(0,1)$.
$\textit{trueWays}=L_T R_F+L_F R_T=1\cdot1+0\cdot0=1$.
$\textit{falseWays}=L_T R_T+L_F R_F=1\cdot0+0\cdot1=0$. So
$f(0,2)=(1,0)$. $f(0,4)$ (the full expression, one operator position
to try, $k=3$, operator \texttt{|}): $L=f(0,2)=(1,0)$, $R=f(4,4)=(0,1)$.
$\textit{trueWays}=L_T R_T+L_T R_F+L_F R_T=1\cdot0+1\cdot1+0\cdot0=1$.
$\textit{falseWays}=L_F R_F=0\cdot1=0$. Final answer:
$\textit{trueWays}=1$ (the single way: $(\texttt{T}\,\texttt{\^{}}\,
\texttt{F})\,\texttt{|}\,\texttt{F}$ evaluates to
$\texttt{True}\,\texttt{|}\,\texttt{False}=\texttt{True}$).

\begin{complexitybox}
\textbf{Time Complexity.} $O(n^3)$ (the same interval-DP bound, with
the operator-only stride halving the constant but not the asymptotic
order).
\textbf{Space Complexity.} $O(n^2)$ for the memo table (each cell
storing a pair).
\end{complexitybox}

\begin{takeawaybox}
Whenever an interval DP's combination step has more than one possible
\emph{outcome} (here, \texttt{True} or \texttt{False}), and a later
combination needs to know the count for \emph{each} outcome to compute
its own result, track a small tuple of values per cell instead of a
single number -- the same upgrade Number of LIS made to its
length-only table.
\end{takeawaybox}
